\section{Related Work}
\label{sec:related}

\subsection{Reinforcement learning for inventory and warehouse robots}
Inventory control is a long-standing test bed for sequential decision making. Gijsbrechts et al.\ showed that a tuned A3C agent can match strong heuristics and approximate dynamic programming on lost-sales, dual-sourcing and multi-echelon problems~\cite{gijsbrechts2022}, and Boute et al.\ survey the design choices such agents require~\cite{boute2022}. Liu et al.\ coordinate the echelons of a supply chain with heterogeneous-agent PPO~\cite{liu2025}. Closer to the store floor, Jullien et al.\ built RetaiL, a grocery restocking environment from real sales data, and trained a distributional DQN to reduce waste~\cite{jullien2023}. On the robotics side, Zhu et al.\ schedule robots in a mobile fulfilment system with deep RL and metaheuristics~\cite{zhu2024}, and Li et al.\ combine PPO with Dijkstra planning for warehouse navigation~\cite{li2024}. These studies decide \emph{how much} to order or \emph{how} to move; the present task couples the two, since the robot's route determines which shelf can be refilled before it runs out. None of them can say how far their policies are from optimal, because their problems are too large to solve exactly. This study deliberately chooses a store small enough to solve, so that every learned policy can be measured against the true optimum.

\subsection{Reward shaping}
Shaping adds reward for progress toward a goal to speed up learning. Randl{\o}v and Alstr{\o}m found that rewarding progress without penalising regress let a bicycle agent profit from riding in circles~\cite{randlov1998}, a simple instance of what Amodei et al.\ call reward hacking~\cite{amodei2016}. Ng, Harada and Russell proved that shaping of the form $\gamma\Phi(s')-\Phi(s)$ preserves the optimal policy, and that this form is also necessary for the guarantee to hold in every MDP~\cite{ng1999}. Wiewiora then showed that, for tabular TD learners, adding such a term is equivalent to initialising $Q(s,a)$ to $\Phi(s)$~\cite{wiewiora2003}. That equivalence turns out to explain most of what shaping does, and does not do, for the tabular learners in Section~\ref{sec:results}.

\subsection{Evaluation practice}
Henderson et al.\ showed that deep RL results often fail to reproduce across random seeds and implementations~\cite{henderson2018}, and Agarwal et al.\ argued for interval estimates over many runs in place of point estimates~\cite{agarwal2021}. Pardo et al.\ showed that treating a time limit as termination corrupts value estimates and recommended bootstrapping at the cut~\cite{pardo2018}. Van Hasselt traced Q-Learning's weakness in stochastic problems to the overestimation introduced by the max operator~\cite{vanhasselt2010}. All four points apply directly to the agents compared here, and Section~\ref{sec:learning} adopts the corresponding practices.
